\documentclass[11pt]{article}
\usepackage[margin=1.1in]{geometry}
\usepackage{amsmath,amssymb,amsthm}
\usepackage{booktabs}
\usepackage{array}
\usepackage[colorlinks=true,allcolors=blue]{hyperref}

\newtheorem{theorem}{Theorem}
\newtheorem{proposition}{Proposition}
\theoremstyle{remark}
\newtheorem{remark}{Remark}

\newcommand{\gr}[1]{\textbf{[#1]}} % epistemic grade tag

\title{A Smith Chart for the Riemann Hypothesis:\\
\large The Cayley--Li Disk as a Unified Instrument Panel for Circle-Side Positivity\\[4pt]
\normalsize IMM Open Series --- companion note to Paper 14 (number to be assigned)}
\author{Travis Bergen\thanks{Independent researcher. ORCID 0009-0009-3950-8390.
Prepared with agent assistance; all numerical claims reproducible from the committed
scripts described in Appendix~\ref{app:repro}.}}
\date{July 15, 2026 --- v0.2 draft (adds $\Theta_\omega$ meter row and E-LI-1b results)}

\begin{document}
\maketitle

\begin{abstract}
\noindent This note proves nothing about the Riemann Hypothesis and presents no evidence for
or against it. Its subject is \emph{instrumentation}: we assemble, under one conformal
coordinate system, the known ``circle-side'' positivity encodings of RH, in exact analogy
with the Smith chart of radio-frequency engineering. The Smith chart is the M\"obius map
$\Gamma=(Z-1)/(Z+1)$ folding the right half-plane of positive-real impedances into the unit
disk, so that passivity becomes legible as geometry. Li's map $z=1-1/s$ performs the same
fold on the critical strip: RH is exactly the statement that every nontrivial zero's image
lies on the unit circle, and the Li coefficients $\lambda_n$ \gr{T: Li 1997;
Bombieri--Lagarias 1999} and the Toeplitz truncations $A(m)$ of the truncated Weil form
\gr{T: Paper 14} are the boundary gauges that would register a departure. Every entry of the
resulting dictionary is graded; the underlying mathematics (M\"obius geometry, Herglotz/
positive-real functions, de Branges canonical systems) is standard and credited, and the
contribution claimed here is assembly and a working, tested instrument only. The instrument
cross-validates against the pre-registered E-LI-1 run: a synthetic off-line quadruple at
$\beta=0.1$, $\gamma\approx98.83$ drives $\mu_{\min}(T_M)<0$ first at $M=272$ in both
implementations.
\end{abstract}

\section{The move, and why engineers made it first}

A two-terminal electrical network is \emph{passive} iff its driving-point impedance $Z(s)$
is a positive-real (Brune) function: analytic on the open right half-plane, real on the real
axis, with $\Re Z(s)\ge 0$ for $\Re s>0$ \gr{T: Brune 1931}. The condition ``a function of a
complex variable keeps a sign on an unbounded region'' is hostile to the eye and to finite
measurement. Smith's 1939 solution was a change of coordinates: the M\"obius map
\begin{equation}
\Gamma \;=\; \frac{Z-1}{Z+1}
\end{equation}
sends the right half-plane bijectively onto the unit disk, the imaginary axis (lossless
impedances) onto the unit circle, and the families $\{\Re Z=\text{const}\}$,
$\{\Im Z=\text{const}\}$ onto two orthogonal families of circular arcs --- the familiar
Smith grid. Positivity becomes containment: \emph{passive $\Leftrightarrow$ inside the disk;
lossless $\Leftrightarrow$ on the rim}. Nothing is proved by the chart; everything becomes
readable on it.

The observation of this note is that the Riemann Hypothesis already possesses the identical
compactification, scattered across three literatures, and that gathering it into one
coordinate system with working gauges is cheap and useful.

\section{The RH chart}

\subsection{The map}

Let $\xi$ be the completed zeta function and $\rho$ range over its nontrivial zeros. Li's
criterion arises from the M\"obius change of variable
\begin{equation}
z \;=\; 1-\frac1s, \qquad s=\frac{1}{1-z},
\label{eq:limap}
\end{equation}
which maps $s=\tfrac12$ to $z=-1$, $s=1$ to $z=0$, and $s=\infty$ to $z=1$.

\begin{proposition}[chart geometry; \gr{T}]
Under \eqref{eq:limap}, $|z| = |s-1|/|s|$. Hence the critical line $\Re s=\tfrac12$ maps
onto the unit circle $|z|=1$ (minus the point $z=1$); the half-plane $\Re s>\tfrac12$ maps
onto the \emph{interior} of the disk; $\Re s<\tfrac12$ onto the exterior. The families
$\{\Re s=\sigma\}$ and $\{\Im s=\gamma\}$ map to circles through $z=1$ and orthogonal arcs
respectively --- a Smith grid for the critical strip.
\end{proposition}

\begin{proof}
$|z|=|1-1/s|=|s-1|/|s|$, and $|s-1|\lessgtr|s|$ according as $s$ is closer to $1$ than to
$0$, i.e.\ according as $\Re s \gtrless \tfrac12$. M\"obius maps send lines to circles
through the image of $\infty$, here $z=1$. \emph{Orientation is pinned by unit test}
\texttt{test\_li\_map\_orientation}; the naive guess (right half maps outside) is wrong,
and the first draft of that test failed for exactly that reason --- recorded here as a
worked example of the verification discipline.
\end{proof}

RH, in chart coordinates \gr{A: restatement, no new content}: \emph{every zero is lossless}
--- the image multiset $\{1-1/\rho\}$ lies on the rim, none strictly inside or outside.

\subsection{The gauges are already circle-side objects}

The point of the assembly is that both standing positivity encodings of RH in this program
are statements \emph{about the chart}, not merely translatable to it.

\textbf{Li gauge.} With $z_\rho:=1-1/\rho$ the Li coefficients are
\begin{equation}
\lambda_n=\sum_\rho\bigl(1-z_\rho^{\,n}\bigr),
\end{equation}
i.e.\ the $n$-th power sums of the charted zeros, and RH $\Leftrightarrow \lambda_n\ge0$ for
all $n\ge1$ \gr{T: Li 1997; Bombieri--Lagarias 1999}. The base of Li's power sum \emph{is}
the charted zero: the chart is not an illustration of the Li criterion, it is its state
space. Under RH, $\lambda_n\sim\tfrac n2\log n$ \gr{T: Lagarias}, which supplies the
normalized margin $s_n=\lambda_n/(\tfrac n2\log(n+1))$ displayed by the instrument.

\textbf{Toeplitz gauge.} Paper 14's sequence $A(m)$ (box basis at scale $\delta$, tent
transform $B$) gives the truncated Weil form as a Toeplitz quadratic form; on the zero side
$A_\zeta(m)=\sum_{\gamma>0}2B(\gamma)\cos(m\delta\gamma)$, and RH $\Leftrightarrow$
$T_M=[A(i-j)]\succeq0$ for all $M,\delta$ \gr{T-conditional: Paper 14, at that paper's own
earned grade}. Positive-definite Toeplitz sequences are exactly the Fourier coefficients of
positive measures on the circle (Herglotz--Carath\'eodory--Toeplitz) \gr{T: classical} ---
again intrinsically a circle-side object.

An off-line zero $\tfrac12\pm\beta\pm i\gamma$ leaves the rim radially
($|z|-1\ne0$, displacement computable in closed form) and enters the gauges as a
$\cosh(m\delta\beta)$-growing term (Toeplitz) and a $(1-z^n)$ defect (Li) --- the two
meters read the same geometric departure at different sensitivities, quantified in
E-LI-1.

\subsection{The impedance column is also already present}

The analogy runs deeper than the disk. The mathematical identity of a positive-real
impedance is a Herglotz/Nevanlinna function, and the RH literature has its own: the Weyl
$m$-function $m=B/A$ of a de Branges canonical system, meromorphic Herglotz, whose poles are
the zeros of $A$ \gr{T: de Branges theory}. Suzuki constructs, for the family
$\Theta_\omega(z)=\xi(\tfrac12-\omega-iz)/\xi(\tfrac12+\omega-iz)$, explicit canonical
systems (unconditionally for $\omega>1$) and states the target criterion: RH as
\emph{positive semidefiniteness of the family of Hamiltonians} \gr{T: Suzuki,
arXiv:1204.1827, for $\omega>1$; Open for the extension}. In network terms: RH is the
assertion that a specific one-port is passive --- the network synthesis literature
(Brune cycles, Kalman--Yakubovich--Popov positivity) and the canonical-system literature are
studying the same function class from opposite ends \gr{T: both classes equal
Herglotz/positive-real up to rotation; A: any suggestion that synthesis \emph{techniques}
transfer}.

\section{The dictionary, graded}

\begin{center}
\small
\begin{tabular}{p{4.6cm}p{5.6cm}l}
\toprule
Smith chart (RF) & RH chart & grade \\
\midrule
Cayley map $\Gamma=\frac{Z-1}{Z+1}$, RHP $\to$ disk & Li map $z=1-1/s$, $\Re s>\tfrac12\to$ interior & \gr{T} \\[2pt]
positive-real (Brune) impedance & Herglotz function / Weyl $m$-function of a canonical system attached to $\xi$ & \gr{T} \\[2pt]
passivity $\Leftrightarrow|\Gamma|\le1$ & RH $\Leftrightarrow$ all charted zeros satisfy $|z|=1$ & \gr{A}\,(restatement) \\[2pt]
lossless network: interlaced imaginary poles/zeros & Hermite--Biehler structure, $A,B$ real interlaced zeros & \gr{T}\,(HB theory) \\[2pt]
reflection coefficient at a frequency & charted zero $z_\rho=1-1/\rho$ = base of Li's power sum & \gr{T}\,(identity) \\[2pt]
VSWR / return-loss meters & $\lambda_n$, $\mu_{\min}(T_M)$ positivity gauges & \gr{T} \\[2pt]
all-pass filter: $|H(j\omega)|=1$, poles strictly stable & inner function $\Theta_\omega$: unimodular on $\mathbb R$; RH tied to no poles in $\mathbb C^+$ & \gr{T}\,(Blaschke; Suzuki) \\[2pt]
matching: move the point to center & no analogue --- zeros cannot be moved; chart is read-only & \gr{A}\,(disanalogy) \\[2pt]
finitely many readings certify a device & positivity needed at \emph{all} depths/windows; no finite certificate & \gr{Open}\,(= RH) \\
\bottomrule
\end{tabular}
\end{center}

The last row is the honest core: a Smith chart certifies a \emph{given, finite} device
because a physical one-port is determined by finitely many parameters. $\xi$ is not. The
chart makes the ``forever'' structure of RH visible; it does not shorten it. Any use of this
note's language that slides from row 3 (restatement) toward a claim of progress on row 8
exceeds its derivation and is disclaimed in advance.

\section{The instrument}

A self-contained interactive chart (HTML/Plotly, dark-panel instrument idiom) accompanies
this note, together with the generating library. Committed parameters: Odlyzko zero table,
first $100{,}000$ zeros ($T\approx74{,}921$); $\delta=0.05$ (Paper 14 primary); display and
gauge depth $n,M\le445$, the binding $D_{\max}$ measured by E-LI-1's tail-stability gate G2
at this zero count. Displays: (i) the charted strip with Smith grid, first 60 zeros on the
rim; (ii) normalized Li margin $s_n$; (iii) Toeplitz margin $\mu_{\min}(T_M)$, log scale;
(iv) a planted synthetic quadruple $\tfrac12\pm0.1\pm i\,98.831$ (toggle), the honest
detection cell identified by E-LI-1's post-hoc analysis.

\textbf{Numerics} \gr{Numerics; truncated, labeled}. Zero-side truncated $\lambda_1=0.0230736$
against the published $0.0230957\ldots$ \gr{T: Keiper--Li} --- the $9\times10^{-4}$ relative
gap is the omitted tail, deliberately not patched here (the E-LI-1 pipeline carries the
mpmath tail integral where gate-grade accuracy is needed). On the true zero set both gauges
stay positive over the full depth range ($\mu_{\min}(T_{445})=4.84\times10^{-6}$). With the
planted quadruple, $\mu_{\min}(T_M)<0$ first at $M=272$, agreeing exactly with E-LI-1's
independently computed detection depth; the global Li margin does not resolve the same
violation by depth 445, reproducing E-LI-1's measured sensitivity asymmetry (its post-hoc
finding that the global $\lambda_n$ statistic has no single-zero power at this depth).

\section{The gauges under load: the E-LI-1b comparison}

After this note's first draft, the pre-registered experiment E-LI-1b (protocol registered
2026-07-15, thresholds frozen before data; successor to E-LI-1 under its permitted
follow-up clause) put all three meters through a powered planted-violation battery:
12 cells ($\varepsilon\in\{0.1,0.05,0.02,0.01\}\times\gamma^*\in\{14.13,\,98.83,\,999.79\}$),
first $2{,}001{,}052$ zeros (Odlyzko), depth $n,M\le2000$, mass-matched double-pair nulls
(the control E-LI-1 lacked), seeds 6000--6019. All validity gates passed; the reproduction
gate matched Paper 14's \S4 numerics to displayed digits. Registered verdict
\gr{Numerics, pre-registered}: \textbf{Toeplitz wins} --- $11$ of $12$ cells resolved, the
$A(m)$ gauge detected in all $11$ (every detection by genuine indefiniteness,
$\mu_{\min}<0$, normalization-invariant), while the Li gauge --- even repaired to a
differenced statistic with single-zero power --- detected in exactly one cell
($\gamma^*=14.13$, $\varepsilon=0.1$, at $n=1244$ vs.\ Toeplitz's $M=87$). Detection-depth
scaling followed the $\cosh(m\delta\varepsilon)$ mechanism, roughly
$M_{\rm det}\propto 1/(\delta\varepsilon)$ with a $\gamma^{*2}$ height penalty from
$B(\gamma)$.

Two registered predictions failed and are reported as such. (i) The $\Theta_\omega$
detection staircase is \emph{not} ``$\varepsilon>\omega$'': the measured law, exact in
$72/72$ probe configurations, is $\varepsilon-\omega\ge v$ --- the planted pole must sit at
or above the probe line $\Im z=v$; a pole below the probe line leaves $|\Theta_\omega|\le1$
on it. The $\Theta_\omega$ gauge is therefore local in both coordinates, and an
$\omega$-sweep must be paired with a downward $v$-sweep. (ii) The registered Li power
estimate was optimistic by a factor $\sim1.7$ in one cell. Full protocol, deviations log
(four entries, all declared before computation, spot-checks passed), gate outcomes, and
machine-readable data accompany the repository.

Instrument conclusion at earned grade: for violation \emph{detection} at these depths and
$\delta=0.05$, the Toeplitz meter is decisively the more sensitive instrument, measured
twice under two null designs; the Li meter's weakness is a property of its global statistic,
not of the Li encoding as mathematics. None of this bears on the truth of RH.

\section{Prior art and what is and is not new}

Nothing in the mathematics of this note is new. The compactification is Li's (1997), with
the arithmetic form by Bombieri--Lagarias (1999) and large-scale $\lambda_n$ computation by
Ma\'slanka. The positive-real/Herglotz view of $\xi$ through canonical systems is de
Branges's program, made explicit as a positivity-of-Hamiltonians criterion by Suzuki
(arXiv:1204.1827); Lagarias's surveys connect the $m$-function language to zeta directly.
The Toeplitz encoding is Paper 14 of this series. The line-side frontier (Weil positivity on
the id\`ele class group, archimedean part proved) is Connes--Consani, arXiv:2006.13771, with
the current operator-theoretic front in Connes--Consani--Moscovici, arXiv:2511.22755. What
this note contributes: the assembled dictionary at explicit grades, the observation that
both of this program's standing gauges are intrinsically objects of one disk coordinate
system, and a tested, reproducible instrument that renders them together. Assembly is
credited as assembly.

\section{What the chart is for}

Three concrete uses inside this program, none of which touch the truth of RH:
(1) \emph{pedagogy and communication} --- the disk-with-rim picture states RH's content
faster than any half-plane figure, and the planted-violation toggle shows what failure
would look like on real gauges; (2) \emph{instrument selection} --- the chart is the natural
cockpit for E-LI-1-style comparisons, and E-LI-1b's design notes (mass-matched nulls,
deeper $D_{\max}$) can be read directly off its displays; (3) \emph{a common coordinate
system} --- future circle-side encodings (inner functions $\Theta_\omega$, Schur/Carath\'eodory
parametrizations) drop onto the same panel, so that ``which encoding detects sooner''
questions stay commensurable.

\appendix
\section{Reproducibility}\label{app:repro}
Workspace \texttt{rh\_smith/}: \texttt{core.py} (maps, $\lambda_n$, $A(m)$, Toeplitz
$\mu_{\min}$; epistemic labels in docstrings), \texttt{make\_chart\_data.py} (all committed
parameters), \texttt{build\_html.py} + \texttt{chart\_template.html} (instrument),
\texttt{test\_core.py} (9 unit tests, all passing: map geometry and orientation, round-trip,
classical Smith sanity, $\lambda_1$ reproduction envelope, positivity on true zeros,
$B$-tent properties, planted-violation detection at $M\le301$). Zero data:
Odlyzko, \texttt{zeros1} table. Runtime: full data generation $\approx11$\,s on the
reference container.

\begin{thebibliography}{9}\small
\bibitem{Brune} O.~Brune, \emph{Synthesis of a finite two-terminal network whose
driving-point impedance is a prescribed function of frequency}, J. Math. Phys. 10 (1931).
\bibitem{Smith} P.~H.~Smith, \emph{Transmission line calculator}, Electronics 12 (1939).
\bibitem{Li} X.-J.~Li, \emph{The positivity of a sequence of numbers and the Riemann
hypothesis}, J. Number Theory 65 (1997) 325--333.
\bibitem{BL} E.~Bombieri, J.~C.~Lagarias, \emph{Complements to Li's criterion for the
Riemann hypothesis}, J. Number Theory 77 (1999) 274--287.
\bibitem{Maslanka} K.~Ma\'slanka, \emph{Effective method of computing Li's coefficients and
their properties}, arXiv:math/0402168.
\bibitem{Suzuki} M.~Suzuki, \emph{A canonical system of differential equations arising from
the Riemann zeta-function}, arXiv:1204.1827.
\bibitem{CC} A.~Connes, C.~Consani, \emph{Weil positivity and trace formula: the archimedean
place}, arXiv:2006.13771.
\bibitem{P14} T.~Bergen, \emph{The Truncated Weil Form as a Toeplitz Sequence} (IMM Paper
14), Zenodo, doi:10.5281/zenodo.20684022.
\bibitem{ELI1} E-LI-1 protocol and results ledger (pre-registered, 2026-07-14), IMM working
documents.
\end{thebibliography}

\end{document}
